\chapter{Det vrange som innsikt}

\begin{motto}
Eit problem som skifter form når du grip det,\\ er ikkje eit dårleg problem. Det er ei anna sort problem.
\end{motto}

Få omgrep frå designteorien har blitt så misbrukte som «vrange problem». Det vandrar gjennom føredrag og søknader som ei orsaking: prosjektet vart ikkje ferdig, men problemet var jo \emph{wicked}; resultatet let seg ikkje måle, men slik er det no eingong med vrange problem; ingen kan seie om grepet var rett, for vrange problem har ingen rette svar. Brukt slik er omgrepet eit alibi — eit ord som gjer ein dygd av at faget ikkje held det det lovar. Og \textsc{Grunnlaust} har rett i å reagere på den bruken. Men reaksjonen råkar bruken, ikkje kjelda. Går ein attende til der omgrepet kjem frå, til Horst Rittel og Melvin Webber i 1973, finn ein ikkje ei orsaking. Ein finn ei av dei skarpaste og mest haldbare innsiktene faget har levert: at det finst ein heil klasse av problem som er strukturelt ulike dei vitskapen løyser, og at å handsame dei som om dei var like, er den eigentlege feilen. Dette kapitlet hentar omgrepet attende frå alibiet og viser at å ta vrangskapen på alvor ikkje er ei unnviking, men sjølve det modne grepet.

\section{Kva Rittel og Webber faktisk sa}

Les originalen, og det fyrste som slår ein, er kor lite han liknar slagordet. Rittel og Webber skreiv ikkje om design i snever forstand, og slett ikkje om designarar som ville sleppe unna ansvar. Dei skreiv om planlegging — om samfunnsplanlegging, politikk, det å gripe inn i ein by eller eit samfunn — og dei skreiv i ein bestemt historisk situasjon: den optimistiske etterkrigstrua på at samfunnsproblem kunne løysast med same systematiske, vitskaplege metode som hadde sett menneske på månen \autocite{rittel1973}. Det dei gjorde, var å vise kvifor denne trua måtte slå feil. Ikkje fordi planleggjarane var udugelege, og ikkje fordi metoden var dårleg, men fordi problema sjølve hadde ein annan struktur enn dei problema metoden var laga for.

Og strukturen skildra dei presist, ikkje ullent. Eit \emph{tamt} problem — eit problem i matematikken, sjakken, eller den klassiske vitskapen — har ei klar formulering; når du har formulert det, veit du kva som skal til for å løyse det. Det har eit stoppunkt: du veit når du er ferdig. Løysinga er rett eller gal. Du kan prøve henne, forkaste henne, prøve på nytt utan kostnad. Eit \emph{vrangt} problem har ingen av desse eigenskapane. Det lèt seg ikkje formulere endeleg, for formuleringa av problemet og forståinga av løysinga er to sider av same sak — du veit ikkje heilt kva problemet er før du ser kva ei løysing ville innebere. Det har ikkje noko stoppunkt; du sluttar fordi tida, pengane eller tolmodet tek slutt, ikkje fordi problemet er løyst. Løysingane er ikkje rette eller gale, berre betre eller verre, og kven som dømer, avheng av kven det går ut over. Kvart vrangt problem er i grunnen eitt av sitt slag. Og — dette er det avgjerande — kvar freistnad på å løyse det endrar problemet, slik at du ikkje får prøve to gonger frå same utgangspunkt. Det finst inga øving. Kvart trekk er for fullt alvor.

Dette er ikkje vag prosa. Det er ei \emph{strukturell} skildring, like presis på sitt felt som ein definisjon i matematikken er på sitt. Rittel og Webber tel opp eigenskapane, set dei opp mot eigenskapane til tamme problem, og viser punkt for punkt korleis dei skil seg. Den som les dette som ei orsaking, har ikkje lese det. Det er ein diagnose.

Verdt å merke seg er kva slags forfattarar dette var. Rittel underviste i designteori og vitskapsteori; Webber var byplanleggjar. Ingen av dei var motstandarar av rasjonalitet eller systematikk — tvert om kom dei begge frå miljø som hadde sett si lit til den. Artikkelen deira er ikkje eit oppgjer med fornufta, men ei sjølvransaking innanfrå: ei erkjenning, frå folk som hadde prøvt den systematiske metoden på samfunnsproblem, av kvifor han ikkje strekte til. Det gjev diagnosen ei tyngd ho ikkje ville hatt om han kom frå nokon som aldri hadde halde metoden høgt. Dei sa ikkje «metoden er verdlaus»; dei sa «her er ein klasse problem metoden ikkje er laga for, og vi har lært det den harde vegen». Å lese dette som eit ynske om å sleppe krav, er å snu erfaringa deira på hovudet.

\section{Kategorifeilen, ein gong til}

No kan ein sjå kvifor skuldinga om at design vik unna falsifiserbarheit, bommar nett her. \textsc{Grunnlaust} hevdar at eit fag utan falsifiserbare påstandar manglar eit fundament, og peikar på at designdommar ikkje lèt seg etterprøve slik vitskaplege hypotesar gjer. Men Rittel og Webber sin innsikt er nett at dette ikkje er ein mangel ved faget; det er ein eigenskap ved problema. Å krevje at eit vrangt problem skal handsamast med falsifisering, er som å krevje at ein arkitekt skal etterprøve om akkurat denne bygningen, på akkurat denne tomta, for akkurat desse menneska, var den rette — ved å byggje henne på nytt under kontrollerte vilkår med éin variabel endra. Eksperimentet kan ikkje gjerast. Ikkje fordi arkitekten er feig, men fordi problemet er av eit slag der øvingsrunden ikkje finst.

Dette er kategorifeilen i kapittel \textsc{i}, sett frå problema si side i staden for frå kunnskapen si side. Der det fyrste kapitlet viste at kravet om naturvitskapleg grunngjeving er ein kategorifeil mot designaren sin kunnskap, viser Rittel og Webber at det same kravet er ein kategorifeil mot designaren sine problem. Det finst problem som lèt seg løyse ved hypotese, eksperiment og falsifisering, og det finst problem som ikkje gjer det — ikkje fordi nokon enno ikkje har funne metoden, men fordi sjølve strukturen i problemet stengjer for han. Å late som om alle problem er av det fyrste slaget, er ikkje strenghet. Det er å sjå berre éin sort problem fordi ein berre har eitt sort verktøy.

Og her gjer \textsc{Grunnlaust} det same grepet som modernismen gjorde med Sullivan: les ein karikatur i staden for originalen. Karikaturen seier at «vrangt problem» er eit ord designarar gøymer seg bak. Originalen seier at det finst ein presist skildra klasse av problem som motstår den vitskaplege metoden av strukturelle grunnar, og at den som ikkje ser det, kjem til å bruke feil reiskap og forundre seg over at det ikkje verkar. Den fyrste lesinga gjer omgrepet til ei svakheit. Den andre — den rette — gjer det til ei av faget sine djupaste innsikter. Skilnaden er ikkje velvilje. Skilnaden er om ein har lese teksten.

\section{Simon, og at det vrange har struktur}

Eitt motargument står att, og det er verdt å ta på fullt alvor, fordi det ser ut til å snu innsikta mot seg sjølv. Om vrange problem verkeleg er heilt ustrukturerte, eitkvart av sitt slag, utan stoppunkt og utan rette svar — er ikkje då sjølve det å studere dei umogleg? Vert ikkje «vrangt problem» då eit ord for «det vi ikkje kan seie noko om»? Og er ikkje det nett alibiet om att, berre i finare språkdrakt?

Svaret kjem frå ein som ikkje var Rittel sin meiningsfelle, men tvert om hans store motpol: Herbert Simon. Same året, 1973, skreiv Simon om det han kalla \emph{ill-structured problems} — dårleg strukturerte problem \autocite{simon1973illstructured}. Der Rittel framheva kor radikalt ulike desse problema var frå dei tamme, ville Simon vise det motsette: at sjølv eit dårleg strukturert problem har ein struktur som kan studerast, og at grensa mellom velstrukturert og dårleg strukturert ikkje er ein avgrunn, men ein glidande overgang. Eit problem er sjeldan dårleg strukturert heile vegen ned; det er samansett av delar, og mange av delane er velstrukturerte. Den dugande problemløysaren arbeider seg innover, gjer det vrange handterleg stykke for stykke, hentar inn kunnskap undervegs som omformar eit ope problem til ei rekkje meir avgrensa. Prosessen sjølv har eit mønster, og mønsteret lèt seg skildre.

Det vakre er kva Rittel og Simon \emph{saman} viser, på tvers av usemja si. Rittel viser at det vrange er ekte: ein kan ikkje definere det bort eller løyse det med vitskapleg metode utan å misforstå det. Simon viser at det vrange likevel har form: det er ikkje kaos, men ein studerbar struktur med eigne reglar for korleis ein arbeider i det. Det fyrste hindrar at vrangskapen vert ein lettvint utveg — for om problemet har struktur, kan ein gjere det betre eller verre, og då finst det noko å svare for. Det andre hindrar at vrangskapen vert ein mur — for om strukturen lèt seg studere, finst det noko å lære, å overføre, å byggje vidare på. Mellom dei to forsvinn alibiet heilt. Den som tek vrangskapen på alvor i Rittel og Simon sin forstand, har ikkje sagt «her kan ingen vite noko». Han har sagt «her gjeld andre reglar enn i laboratoriet, og dei reglane kan vi kartleggje».

Dette var alt føresett i \textcite{simon1969sciences} sitt prosjekt om ein vitskap om det kunstige: at det menneskeskapte har ein eigen logikk som er like studerbar som naturen sin, berre annleis, fordi det handlar om korleis ting \emph{kan} bli og ikkje korleis dei \emph{er}. Designproblem er det fremste dømet på denne logikken. Dei er kontingente: dei kunne ha vore annleis, og oppgåva er å gjere dei betre, ikkje å fastslå korleis dei nødvendigvis er.

At Simon og Rittel kom til kvar si vektlegging, er for øvrig ikkje eit teikn på at faget ikkje veit kva det meiner, men på det motsette: at det har ein indre debatt med to klart formulerte posisjonar som kvar fangar noko sant. Simon ville helst gjere det dårleg strukturerte velstrukturert — bryte det ned til det vart handterleg — og frykta at å overdrive vrangskapen ville lamme. Rittel ville halde fast ved at noko i problema motstår all slik nedbryting, og frykta at å undervurdere vrangskapen ville føre til arrogante, skadelege inngrep. Begge fryktene er rettkomne, og ein moden utøvar lever i spennet mellom dei: han bryt ned det som lèt seg bryte ned, og er audmjuk overfor det som ikkje gjer det. Eit fag som rommar den spenninga, og som har tenkt henne grundig nok til å gje begge sider eit namn, er ikkje eit fag utan grunnlag. Det er eit fag som har forstått problema sine godt nok til å vere usamd om dei på ein opplyst måte.

\section{Rittel mot det anti-teoretiske}

Den sterkaste tilbakevisinga av alibi-lesinga er likevel kva Rittel sjølv gjorde etterpå. Om «vrangt problem» verkeleg var ei orsaking for at ein ikkje treng grunngje noko, ville den som fann opp omgrepet, ha brukt det til å sleppe unna argumentasjon. Rittel gjorde det stikk motsette. Han bruka resten av karrieren på å byggje strenge metodar for å argumentere gjennom nett desse problema. Innsikta om at problema var vrange, var for han ikkje slutten på rasjonaliteten, men byrjinga på ein \emph{annan} rasjonalitet — det han kalla ei argumentativ tilnærming til planlegging.

Frukta av dette var system for å føre og prøve resonnement i situasjonar utan rette svar: metodar der ein gjer spørsmåla eksplisitte, listar opp stillingstaka til kvart spørsmål, og legg fram grunnane for og imot kvar posisjon så dei kan vegast mot kvarandre i open dag. Poenget var å gjere det skjulte synleg — å tvinge fram, svart på kvitt, kva som står på spel, kva alternativa er, og kvifor ein vel som ein vel. Dette er det motsette av eit alibi. Eit alibi gøymer; Rittel sin metode avdekkjer. Eit alibi seier «du kan ikkje krevje grunnar av meg»; Rittel sin metode seier «her er grunnane mine, lagde fram så du kan prøve dei». At svaret ikkje kan falsifiserast som ein fysisk hypotese, tyder ikkje at det ikkje kan grunngjevast og prøvast; det tyder at grunngjevinga og prøvinga har ei anna form — argumentativ, ikkje eksperimentell.

Det er dette \textsc{Grunnlaust} ikkje får auge på når omgrepet vert lese som ei unnviking. Mannen som peika på at designproblem er vrange, var ikkje anti-teoretisk. Han var blant dei mest systematisk teoretiske tenkjarane faget har hatt, og han vart det \emph{av} innsikta om vrangskapen, ikkje på trass av henne. Innsikta kravde ein metode, og han bygde han. \textcite{buchanan1992wicked} såg dette klart då han plasserte dei vrange problema i hjartet av designtenkinga: ikkje som eit problem faget har, men som forklaringa på kva slags fag design er — eit fag for det ubestemte, der oppgåva nett er å gje form til det som ikkje har éi rett form.

Då fell den enklaste versjonen av skuldinga. Å seie at eit problem er vrangt, er ikkje å seie at alt er like godt, at ingen kan dømme, eller at faget slepp unna. Det er å seie at vi står overfor ein klasse problem som krev ein eigen, krevjande form for argumentasjon — og at det modne er å utvikle den argumentasjonen, slik Rittel gjorde, ikkje å late som problema er tamme, slik kritikaren gjer når han krev falsifisering. Vrangskapen er ikkje faget si svakheit. Rett forstått er han faget sitt studieobjekt.

\vignett

Men om design har sine eigne problem, treng det òg ein eigen plass i kunnskapen sitt hus. Kva slags disiplin er det som arbeider slik? Det er emnet for neste kapittel.

\vidarelesing{\fullcite{rittel1973}; \fullcite{simon1973illstructured}; \fullcite{buchanan1992wicked}; \fullcite{simon1969sciences}.}
